
\documentclass[11pt]{article}
\usepackage[margin=1in]{geometry}
\usepackage[T1]{fontenc}
\usepackage[utf8]{inputenc}
\usepackage{amsmath,amssymb,amsthm}
\usepackage{booktabs}
\usepackage{array}
\usepackage{enumitem}
\usepackage{xcolor}
\usepackage{hyperref}
\usepackage{url}
\hypersetup{colorlinks=true,linkcolor=blue!50!black,citecolor=blue!50!black,urlcolor=blue!50!black}
\setlist[itemize]{leftmargin=1.4em,itemsep=0.25em,topsep=0.25em}
\setlist[enumerate]{leftmargin=1.6em,itemsep=0.25em,topsep=0.25em}
\newcommand{\WMPLN}{WM--PLN}
\newcommand{\PLN}{PLN}
\newcommand{\KS}{K\&S}
\newcommand{\codepath}[1]{\path{#1}}
\newcommand{\leanref}[1]{\path{#1}}
\newcommand{\Evidence}{\ensuremath{\mathcal{E}}}
\newcommand{\World}{\ensuremath{\mathcal{W}}}
\newcommand{\Query}{\ensuremath{\mathcal{Q}}}
\newcommand{\R}{\ensuremath{\mathbb{R}}}
\newcommand{\draftnote}{\noindent\textit{Draft extended abstract for author revision.}}
\sloppy

% Draft inspiration only. Author should rewrite, verify claims, and choose final voice.
\title{PLN Chaining as Certified World-Model Extraction}
\author{Zarathustra Amareis Goertzel\\\small Mettapedia Project}
\date{Draft for JoeFest extended abstract}
\begin{document}
\maketitle
\draftnote

\begin{abstract}
PLN chaining is attractive because it offers local, fast uncertain inference. It is also dangerous: local truth values do not, in general, determine the joint or causal world model from which they were extracted. This extended abstract proposes a precise status for PLN chaining: it is a family of certified, side-conditioned rewrites compiled from declared world-model classes. In the world-model calculus, a state $W$ is the semantic object, extraction maps $W$ and a query $q$ to evidence, and scalar strength/confidence values are derived views. Classical PLN deduction, induction, and abduction are then recovered as theorems for specific structures such as Bayesian-network chains and forks, while collider and hidden-dependence cases explain why unrestricted chaining fails. The result is a Halpern-style account of PLN: local rules are useful, but only when their model-theoretic preconditions are explicit.
\end{abstract}

\section{The problem with unconstrained chaining}

The original PLN rule family was built to support fast uncertain inference over link truth values \cite{GoertzelPLN2009}. A typical rule combines quantities such as $P(B \mid A)$ and $P(C \mid B)$ to estimate $P(C \mid A)$. Operationally this is valuable. Semantically it is incomplete: the same local truth values may arise from different joint distributions, causal structures, or evidence histories.

This is exactly the kind of representational issue that runs through Halpern's work on uncertainty and causality. A formalism must say not only which numbers are manipulated, but what model makes those manipulations sound \cite{Halpern2003,HalpernPearl2005}. The world-model calculus answers by separating three layers:
\begin{itemize}
\item a \emph{world-model state} $W$, which is the revisable semantic object;
\item an \emph{evidence carrier} extracted from $W$ by queries;
\item \emph{views} such as strength and confidence, used for control and display.
\end{itemize}

A PLN chaining rule is therefore not a primitive law of uncertain logic. It is a compiled shortcut whose validity depends on the model class and its side conditions.

\section{Exact deduction from a chain model}

Consider a Bayesian-network chain $A \to B \to C$. If $C$ is screened off from $A$ by $B$ and by $\neg B$, and $0<P(B)<1$, the standard PLN deduction formula is exact:
\[
  \mathrm{strength}(W, A \to C)
  = \mathrm{plnDed}(s_{AB},s_{BC},p_B,p_C),
\]
where
\[
\begin{gathered}
  s_{AB}=\mathrm{strength}(W,A\to B),\qquad
  s_{BC}=\mathrm{strength}(W,B\to C),\\
  p_B=\mathrm{strength}(W,B),\qquad
  p_C=\mathrm{strength}(W,C).
\end{gathered}
\]
In Lean this is represented by theorem families such as
\leanref{chainBN_plnDeductionStrength_exact} in
\codepath{Mettapedia/Logic/PLNBayesNetFastRules.lean}. A corresponding fork theorem handles $A \leftarrow B \to C$.

The significance is not that the algebraic formula is new. The significance is that the formula has been relocated. It is no longer a free-floating truth-value propagation rule. It is the projection, through the strength view, of a valid world-model extraction theorem.

\section{Negative controls: fork, chain, collider}

The most convincing presentation should include both positive and negative lanes.

\begin{center}
\begin{tabular}{p{0.22\linewidth}p{0.32\linewidth}p{0.34\linewidth}}
\toprule
\textbf{Pattern} & \textbf{Expected behavior} & \textbf{Message} \\
\midrule
Chain $A \to B \to C$ & deduction exact under screening-off & PLN chaining is a certified shortcut \\
Fork $A \leftarrow B \to C$ & related exact rule under common-cause conditions & side conditions expose what is being assumed \\
Collider $A \to B \leftarrow C$ & local chaining blocked or context-dependent & local truth values alone do not determine the joint \\
Hidden confounder & scalar propagation can be misleading & the missing world model matters \\
\bottomrule
\end{tabular}
\end{center}

This is a compact tribute to Halpern's semantic method: do not ban local rules; state the models that license them.

\section{Evidence-carrying chaining}

A second refinement is to propagate evidence, not only scalar strength. Binary evidence carries counts $(n^+,n^-)$, from which strength and confidence are extracted:
\[
  s = \frac{n^+}{n^+ + n^-}, \qquad
  c = \frac{n^+ + n^-}{n^+ + n^- + \kappa}.
\]
If two links have the same strength but different counts, they should not contribute equally to downstream decisions. Evidence-carrying chaining therefore supports principled revision, confidence-aware abstention, and provenance tracking. It also clarifies what ProbLog-style point probabilities discard: the sufficient statistics behind the estimate.

A minimal experiment for the extended abstract is:
\begin{enumerate}
\item instantiate chain, fork, collider, and hidden-confounder networks;
\item compute exact query answers from the world model;
\item compare raw scalar PLN chaining, certified side-conditioned chaining, and evidence-carrying chaining;
\item report when the shortcut is exact, when it is blocked, and when it becomes an interval or abstention result.
\end{enumerate}

\section{Conclusion}

This submission would frame WM--PLN as a disciplined answer to a long-standing PLN question: when does chaining actually work? The answer is not rhetorical. It works when it is a theorem about a model class. It fails, weakens, or abstains when the required semantic information is absent. That is a cleaner, more honest, and more Halpernian role for PLN rules.


\begin{thebibliography}{99}

\bibitem{Halpern2003}
J.~Y. Halpern.
\newblock \emph{Reasoning about Uncertainty}.
\newblock MIT Press, 2003.

\bibitem{FaginHalpernMegiddo1990}
R.~Fagin, J.~Y. Halpern, and N.~Megiddo.
\newblock A logic for reasoning about probabilities.
\newblock \emph{Information and Computation}, 87(1--2):78--128, 1990.

\bibitem{FaginHalpernMosesVardi1995}
R.~Fagin, J.~Y. Halpern, Y.~Moses, and M.~Y. Vardi.
\newblock \emph{Reasoning about Knowledge}.
\newblock MIT Press, 1995.

\bibitem{Halpern1999}
J.~Y. Halpern.
\newblock Cox's theorem revisited.
\newblock \emph{Journal of Artificial Intelligence Research}, 11:429--435, 1999.

\bibitem{HalpernPearl2005}
J.~Y. Halpern and J.~Pearl.
\newblock Causes and explanations: A structural-model approach. Parts I and II.
\newblock \emph{British Journal for the Philosophy of Science}, 56(4):843--887 and 56(4):889--911, 2005.

\bibitem{GoertzelPLN2009}
B.~Goertzel, M.~Ikl\'{e}, I.~F. Goertzel, and A.~Heljakka.
\newblock \emph{Probabilistic Logic Networks: A Comprehensive Framework for Uncertain Inference}.
\newblock Springer, 2009.

\bibitem{KnuthSkilling2012}
K.~H. Knuth and J.~Skilling.
\newblock Foundations of inference.
\newblock \emph{Axioms}, 1(1):38--73, 2012.

\bibitem{Pearl2000}
J.~Pearl.
\newblock \emph{Causality: Models, Reasoning, and Inference}.
\newblock Cambridge University Press, 2000.

\end{thebibliography}

\end{document}
